% PUC-Rio 1999/2001/2002 — 19 questões (12 MCQ + 7 discursivas)

---
tipo: Discursiva
dificuldade: Média
nivel: médio
disciplina: Matemática
assunto: Sistemas Lineares
resposta: Para o sistema ter infinitas soluções, as equações devem ser proporcionais: $\dfrac{a}{2} = \dfrac{5}{3} = \dfrac{b}{6}$, portanto $a = \dfrac{10}{3}$ e $b = 10$.
tags: [sistema dependente, infinitas soluções, proporcionalidade]
fonte: concurso
concurso: PUC
banca: PUC
ano: 1999
numero: 1
---
\question
Ache os valores de $a$ e $b$ para que o sistema
$$\begin{cases} 2x + 3y = 6 \\ ax + 5y = b \end{cases}$$
tenha mais do que uma solução.

---
tipo: Discursiva
dificuldade: Média
nivel: médio
disciplina: Matemática
assunto: Trigonometria
resposta: $\operatorname{sen}(x) + \cos(x) = 0 \Rightarrow \operatorname{tg}(x) = -1$. No intervalo $[0, 2\pi]$, a equação $\operatorname{tg}(x) = -1$ possui exatamente 2 soluções: $x = \dfrac{3\pi}{4}$ e $x = \dfrac{7\pi}{4}$.
tags: [equação trigonométrica, tangente, intervalo]
fonte: concurso
concurso: PUC
banca: PUC
ano: 1999
numero: 2
---
\question
Quantas soluções de $\operatorname{sen}(x) + \cos(x) = 0$ existem para $x$ entre $0$ e $2\pi$?

---
tipo: Discursiva
dificuldade: Média
nivel: médio
disciplina: Matemática
assunto: Equação do Primeiro Grau
resposta: Sejam os sete inteiros consecutivos $n, n+1, n+2, n+3, n+4, n+5, n+6$. A soma dos quatro primeiros é $4n+6$ e a dos três últimos é $3n+15$. Igualando: $n = 9$. Os números são $9, 10, 11, 12, 13, 14, 15$.
tags: [inteiros consecutivos, soma, sequência]
fonte: concurso
concurso: PUC
banca: PUC
ano: 1999
numero: 3
---
\question
Ache sete números inteiros consecutivos tais que a soma dos primeiros quatro seja igual à soma dos últimos três.

---
tipo: Discursiva
dificuldade: Média
nivel: médio
disciplina: Matemática
assunto: Princípio Fundamental da Contagem
resposta: Com 12 meses e no máximo 3 pessoas por mês, o grupo pode ter até $3 \times 12 = 36$ pessoas sem que 4 nasçam no mesmo mês. Portanto, o menor tamanho que garante 4 pessoas no mesmo mês é $37$.
tags: [princípio das gavetas, meses, grupo]
fonte: concurso
concurso: PUC
banca: PUC
ano: 1999
numero: 4
---
\question
Qual é o menor número de pessoas num grupo para garantir que pelo menos 4 pessoas nasceram no mesmo mês?

---
tipo: Múltipla Escolha
dificuldade: Fácil
nivel: médio
disciplina: Matemática
assunto: Sistemas Lineares
gabarito: D
tags: [infinitas soluções, sistema linear, 3 variáveis]
fonte: concurso
concurso: PUC
banca: PUC
ano: 2002
numero: 5
---
\question
Assinale a afirmativa correta. O sistema
$$\begin{cases} x + y + z = 1 \\ x + y - z = 1 \end{cases}$$
\begin{choices}
\choice não tem solução.
\choice tem uma solução única $x = 1$, $y = 0$, $z = 0$.
\choice tem exatamente duas soluções.
\CorrectChoice tem uma infinidade de soluções.
\choice tem uma solução com $z = 1$.
\end{choices}

---
tipo: Múltipla Escolha
dificuldade: Fácil
nivel: médio
disciplina: Matemática
assunto: Combinações
gabarito: B
tags: [campeonato, combinações, jogos]
fonte: concurso
concurso: PUC
banca: PUC
ano: 2002
numero: 6
---
\question
O campeonato brasileiro tem, em sua primeira fase, 28 times que jogam todos entre si. Nesta etapa, o número de jogos é de:
\begin{choices}
\choice $376$
\CorrectChoice $378$
\choice $380$
\choice $388$
\choice $396$
\end{choices}

---
tipo: Múltipla Escolha
dificuldade: Média
nivel: médio
disciplina: Matemática
assunto: Princípio Fundamental da Contagem
gabarito: C
tags: [senha, caracteres, alfabético, numérico]
fonte: concurso
concurso: PUC
banca: PUC
ano: 2002
numero: 7
---
\question
A senha de acesso a um jogo de computador consiste em quatro caracteres alfabéticos ou numéricos, sendo o primeiro necessariamente alfabético. O número de senhas possíveis será:
\begin{choices}
\choice $36^4$
\choice $10 \times 36^3$
\CorrectChoice $26 \times 36^3$
\choice $26^4$
\choice $10 \times 26^4$
\end{choices}

---
tipo: Múltipla Escolha
dificuldade: Fácil
nivel: médio
disciplina: Matemática
assunto: Juros Compostos
gabarito: C
tags: [inflação, juros compostos, acumulado]
fonte: concurso
concurso: PUC
banca: PUC
ano: 2002
numero: 8
---
\question
Uma inflação mensal de $2\%$ acumula durante quatro meses uma inflação de, aproximadamente:
\begin{choices}
\choice $7\%$
\choice $9\%$
\CorrectChoice $8{,}25\%$
\choice $10\%$
\choice $12\%$
\end{choices}

---
tipo: Múltipla Escolha
dificuldade: Média
nivel: médio
disciplina: Matemática
assunto: Probabilidade Simples
gabarito: A
tags: [comissão, probabilidade, combinações]
fonte: concurso
concurso: PUC
banca: PUC
ano: 2002
numero: 9
---
\question
De uma turma de 30 alunos é escolhida uma comissão de 3 representantes. Qual a probabilidade de você fazer parte da comissão?
\begin{choices}
\CorrectChoice $\dfrac{1}{10}$
\choice $\dfrac{1}{12}$
\choice $\dfrac{5}{24}$
\choice $\dfrac{1}{3}$
\choice $\dfrac{2}{9}$
\end{choices}

---
tipo: Múltipla Escolha
dificuldade: Média
nivel: médio
disciplina: Matemática
assunto: Equação do Segundo Grau
gabarito: C
tags: [discriminante, raízes reais, raiz dupla]
fonte: concurso
concurso: PUC
banca: PUC
ano: 2002
numero: 10
---
\question
Assinale a afirmativa correta. O polinômio $x^2 - ax + 1$:
\begin{choices}
\choice tem sempre duas raízes reais.
\choice tem sempre uma raiz real.
\CorrectChoice tem exatamente uma raiz real para $a = 2$.
\choice tem exatamente uma raiz real para infinitos valores de $a$.
\choice tem exatamente uma raiz real para $a = 0$.
\end{choices}

---
tipo: Múltipla Escolha
dificuldade: Média
nivel: médio
disciplina: Matemática
assunto: Inequações Modulares
gabarito: D
tags: [módulo, inequação, valor absoluto]
fonte: concurso
concurso: PUC
banca: PUC
ano: 2002
numero: 11
---
\question
Assinale a afirmativa correta. A inequação $-|x| < x$:
\begin{choices}
\choice nunca é satisfeita.
\choice é satisfeita em $x = 0$.
\choice é satisfeita para $x$ negativo.
\CorrectChoice é satisfeita para $x$ positivo.
\choice é sempre satisfeita.
\end{choices}

---
tipo: Múltipla Escolha
dificuldade: Difícil
nivel: médio
disciplina: Matemática
assunto: Progressão Aritmética
gabarito: C
tags: [progressão geométrica, idades, AM-GM, impossibilidade]
fonte: concurso
concurso: PUC
banca: PUC
ano: 2002
numero: 12
---
\question
Um senhor tem $a$ anos de idade, seu filho tem $f$ anos e seu neto, $n$. Sobre estes valores, podemos afirmar:
\begin{choices}
\choice É impossível que $a$, $f$ e $n$ estejam em progressão aritmética.
\choice É impossível que $a$, $f$ e $n$ estejam em progressão geométrica.
\CorrectChoice É impossível que $a$, $f$ e $n$ estejam simultaneamente em progressão aritmética e geométrica.
\choice É possível que $a$, $f$ e $n$ estejam simultaneamente em progressão aritmética e geométrica.
\choice É possível que $a$, $f$ e $n$ estejam em progressão aritmética, mas é impossível que estejam em progressão geométrica.
\end{choices}

---
tipo: Múltipla Escolha
dificuldade: Fácil
nivel: médio
disciplina: Matemática
assunto: Médias
gabarito: E
tags: [média ponderada, pesos, notas]
fonte: concurso
concurso: PUC
banca: PUC
ano: 2002
numero: 13
---
\question
Um aluno faz 3 provas com pesos 2, 2 e 3. Se ele tirou 2 e 7 nas duas primeiras, quanto precisa tirar na terceira para ficar com média maior ou igual a 6?
\begin{choices}
\choice Pelo menos $4$.
\choice Pelo menos $5$.
\choice Pelo menos $6$.
\choice Pelo menos $7$.
\CorrectChoice Pelo menos $8$.
\end{choices}

---
tipo: Múltipla Escolha
dificuldade: Média
nivel: médio
disciplina: Matemática
assunto: Teorema de Pitágoras
gabarito: C
tags: [retângulo, diagonal, área, sistema]
fonte: concurso
concurso: PUC
banca: PUC
ano: 2002
numero: 14
---
\question
Se um retângulo tem diagonal medindo $10$ e lados cujas medidas somam $14$, qual sua área?
\begin{choices}
\choice $24$
\choice $32$
\CorrectChoice $48$
\choice $54$
\choice $72$
\end{choices}

---
tipo: Discursiva
dificuldade: Média
nivel: médio
disciplina: Matemática
assunto: Inequação do Primeiro Grau
resposta: Sendo $F$ o faturamento, o lucro é $F - 0{,}30F - 100{.}000 \geq 40{.}000$, ou seja, $0{,}70F \geq 140{.}000$, portanto $F \geq 200{.}000$.
tags: [lucro, faturamento, imposto, inequação]
fonte: concurso
concurso: PUC
banca: PUC
ano: 2002
numero: 15
---
\question
Uma indústria opera com custo fixo de R\$\,100.000 por ano e paga impostos sobre 30\% de seu faturamento bruto. Quanto deve faturar para que seu lucro no ano seja de no mínimo R\$\,40.000?

---
tipo: Discursiva
dificuldade: Média
nivel: médio
disciplina: Matemática
assunto1: Polinômios
tags1: [avaliação, substituição]
resposta1: $p(0) = -1$; $p(1) = 2$; $p(-1) = 0$; $p(2) = 15$; $p(-2) = -1$
assunto2: Polinômios
tags2: [raízes, fatoração, Bhaskara]
resposta2: $x = -1$; $x = \dfrac{-1+\sqrt{5}}{2}$; $x = \dfrac{-1-\sqrt{5}}{2}$
fonte: concurso
concurso: PUC
banca: PUC
ano: 2002
numero: 16
---
\question
Considere o polinômio $p(x) = x^3 + 2x^2 - 1$.
\begin{parts}
\part Calcule $p(0)$, $p(1)$, $p(-1)$, $p(2)$ e $p(-2)$.
\part Ache as três soluções da equação $x^3 + 2x^2 = 1$.
\end{parts}

---
tipo: Discursiva
dificuldade: Média
nivel: médio
disciplina: Matemática
assunto: Teorema de Pitágoras
resposta: As soluções inteiras são $(a, b) = (1, 6)$ e $(a, b) = (3, 2)$. Em ambos os casos a área é $\dfrac{1}{2} \cdot b \cdot a = 3$.
tags: [triângulo isósceles, inteiros, área]
fonte: concurso
concurso: PUC
banca: PUC
ano: 2002
numero: 17
---
\question
Seja $T$ um triângulo isósceles de base $b$ e altura $a$, onde $a$ e $b$ são inteiros. Dado que os lados iguais de $T$ medem $\sqrt{10}$, calcule a área de $T$.

---
tipo: Múltipla Escolha
dificuldade: Média
nivel: médio
disciplina: Matemática
assunto: Inscrição e Circunscrição
gabarito: A
tags: [triângulo equilátero, circunscrito, inscrito, razão de raios]
fonte: concurso
concurso: PUC
banca: PUC
ano: 2001
numero: 18
---
\question
Qual a razão entre os raios dos círculos circunscrito e inscrito de um triângulo equilátero de lado $a$?
\begin{choices}
\CorrectChoice $2$
\choice $\sqrt{3}$
\choice $\sqrt{2}$
\choice $3a$
\choice $\sqrt{3a^2}$
\end{choices}

---
tipo: Múltipla Escolha
dificuldade: Média
nivel: médio
disciplina: Matemática
assunto: Sistemas Lineares
gabarito: D
tags: [sistema homogêneo, espaço solução, múltiplos]
fonte: concurso
concurso: PUC
banca: PUC
ano: 2001
numero: 19
---
\question
O conjunto de todas as soluções do sistema
$$\begin{cases} x + 2y + 3z = 0 \\ 4x + 5y + 6z = 0 \end{cases}$$
\begin{choices}
\choice é vazio.
\choice consiste apenas no vetor nulo $(0, 0, 0)$.
\choice consiste apenas no vetor $(1, -2, 1)$.
\CorrectChoice consiste em todos os múltiplos de $(1, -2, 1)$.
\choice consiste em todos os múltiplos de $(1, 1, -2)$.
\end{choices}
